% Ch.5 Methodology (~8-10 pp). Draft on Day 3.
\chapter{Methodology}
\label{ch:methodology}

\section{Factor Construction Pipeline}
\label{sec:meth-pipeline}
% Trend-inflation EWMA -> cycle regressions -> local CF -> global GCF (GDP
% weights) -> FX-adjusted FXGCF. Map each step to the code in
% data preperation.R; state the maturity menu {1,2,4,5,9,10}.

\section{Predictive-Regression Design}
\label{sec:meth-design}
% 12-month overlapping returns; duration standardization; the Eq (18)-(23)
% specifications and what each tests.

\section{Statistical Inference}
\label{sec:meth-inference}
% Newey-West HAC \citep{neweywest1987} t-stats and the reverse-regression
% delta method \citep{hodrick1992, weiwright2013}: explain why CP use it and the
% honest caveat (we report HAC; exact Wei-Wright needs monthly returns + the
% appendix -- see cp_inference.R header). BIC relative probabilities; joint Wald
% tests; stationary block bootstrap \citep{politis1994} for small-sample bands
% and R^2 confidence intervals.

\section{Expectations-Hypothesis Monte Carlo}
\label{sec:meth-ehmc}
% Section-1 two-factor EH simulation (cp_montecarlo.R): null distribution of the
% predictive R^2. State the calibration and the documented level gap vs CP.

\section{Out-of-Sample Evaluation}
\label{sec:meth-oos}
% Fully-recursive factors (oos.R) and the Campbell-Thompson \citep{campbellthompson2008}
% out-of-sample R^2 vs the recursive-mean benchmark.

\section{Deviations from the Original Studies}
\label{sec:meth-deviations}
% Be explicit: maturity menu, sample span, HAC-vs-Wei-Wright, 3-maturity
% individual returns. These frame the replication chapter's comparisons.
